\section{Related Work}
% This section first reviews the related work in general mutlimodal pretraining with a focus on contrastive pretraining, and further reviews the progress of Chinese multimodal pretraining. 

% \paragraph{Multimodal Pretraining}
% There are mainly two lines of studies in multimodal pretraining, namely generative multimodal pretraining and contrastive multimodal pretraining. 
% The core of generative multimodal pretraining is the transfer of BERT or T5 to cross-modal data, and a number of models consecutively created new state-of-the art performance in a series of cross-modal downstream tasks~\citep{uniter, unicoder-vl, visualbert, vilbert, interbert, oscar, pixelbert, e2e-vlp, vinvl, clipvil, simvlm, ofa, albef, blip, mplug, vlmo, beit3}, e.g., image captioning~\citep{coco_cap}, VQA~\citep{vqa, vqav2}, etc. 
Previous vision-language pretrained models are mostly BERT/T5-style~\citep{bert, t5}, which involves cross-modal fusion~\citep{uniter, unicoder-vl, visualbert, vilbert, interbert, oscar, pixelbert, e2e-vlp, vinvl, clipvil, simvlm, ofa, albef, blip, mplug, vlmo, beit3}. 
CLIP~\citep{clip}, instead, is a contrastive-learning-based two-tower model, which can serve as a vision foundation model.  
% CLIP demonstrated that the simple contrastive learning with large-scale multimodal data could build an effective vision foundation model. 
% It revolutionized both computer vision and multimodal representation learning. 
Following CLIP, a series of similar contrastive-learning-based multimodal pretrained models were proposed and reached new SOTAs in cross-modal retrieval and zero-shot classification~\citep{align, filip, florence}. 
Furthermore, CLIP can be adaptive to other models. 
A typical case is that CLIP is essential to many image generation models, e.g., DALL-E~\citep{dalle}, DALL-E 2~\citep{dalle2}, Stable Diffusion~\citep{latent_diffusion}, etc. 
% \paragraph{Chinese Multimodal Pretraining}
The success of multimodal pretraining encouraged the transfer of the existing methods to Chinese pretraining, including generative pretrained models~\citep{m6, wenlan, m6-t, m6-10t, ofa} and contrastive pretrained models~\citep{wenlan, wukong, r2d2, altclip}. 
% Note that in this work, we compare Chinese CLIP with the contrastive-learning-based models, 

% Similarly, researchers have proposed models for both generative and contrastive pretraining. \citet{m6} proposed a model that unifies understanding and generation and pretrained on an in-house large-scale dataset, and further extended it to larger scales~\citep{m6, m6-10t}. 
% \citet{wenlan} started the research in contrastive multimodal pretraining for Chinese data and incorporated complex design for better performance. \citet{wukong} proposed the Wukong dataset, a large-scale public dataset for Chinese multimodal pretraining, and the corresponding pretrained model. 
% \citet{r2d2} proposed a new dataset called ZERO whose subset was released, and pretrained a new model called R2D2 that incorporated contrastive two-tower architecture and interactive module for cross-modal fusion. 
% In this work, different from the previous studies that explored intricate model designs, we build a baseline Chinese CLIP without significant architectural changes from the original CLIP, and we use the simple two-stage pretraining to help Chinese CLIP reach outstanding performance. 